\chapter{Discussion}
\label{ch:discussion}

The discussion below examines how the proposed approach answers the research questions posed in Section~\ref{sec:research_questions}, drawing on the quantitative and qualitative results of Chapter~\ref{ch:evaluation}, and addresses the limitations that qualify each answer.

\section{Descriptor Design}
\label{sec:disc_descriptors}
The first research question asked which guitar-specific symbolic descriptors best capture the technical difficulty of classical guitar pieces. The audits of Section~\ref{sec:descriptors} and the qualitative analysis of Section~\ref{sec:qualitative} give a consistent answer: piece length, fretboard coverage and position statistics (\texttt{fret\_entropy}, \texttt{high\_position\_ratio}, \texttt{avg\_fret}), shift frequency, and --- once rhythm is available --- demand-per-beat measures carry the strongest signal, with open-string content as the strongest negative indicator. Equally informative is what did \emph{not} work: several intuitively plausible descriptors, including raw shift distance, string-jump statistics, and polyphony averages, proved nearly uncorrelated with graded difficulty until they were placed in temporal context. The comparison of Section~\ref{sec:context_aware} shows why: a physical demand only becomes difficult in proportion to the time available to execute it, so context-free averages conflate leisurely and virtuosic instances of the same motion.

These findings both extend and depart from the prior work on guitar playability. The descriptors of V{\'a}squez et al.~\cite{vasquez2023quantifying} --- barre usage, chord-shape displacement --- have direct counterparts in our set, but the strongest descriptors for classical repertoire turned out to be ones with no analogue in strummed chord charts: fretboard coverage entropy and rhythm-normalized movement statistics, reflecting the horizontal, position-based way the classical left hand traverses the neck.

Three limitations qualify the descriptor layer. First, tablature commits each note to a string and fret, which the descriptors take at face value; where a transcription chose an awkward fingering, the descriptors measure the transcription, not the piece's best playable realization. A fingering-optimization step (Section~\ref{sec:future_work}) would decouple the two. Second, the rhythm-alignment pipeline collapses multi-voice guitar textures into a single stream of note attacks: pitch content and onset positions --- what the descriptors actually consume --- are correct, but the notated voice-leading structure is lost, so descriptors that would distinguish a sustained bass line under a melody from block writing are out of reach for the PDF-derived majority of the dataset. Third, rhythm recovery requires a matched MIDI performance per PDF score; for 136 pieces the alignment coverage was too low and for 3 no MIDI existed, and those pieces fall back to median-imputed rhythm descriptors. As built, the system cannot grade an arbitrary new PDF score without either native rhythm information or a human-in-the-loop MIDI match.

\section{Predictive Performance}
\label{sec:disc_performance}
The second research question asked whether an interpretable, rubric-based model can predict piece-level difficulty with useful accuracy. The results of Section~\ref{sec:quantitative} support an affirmative answer with a precise qualification. On raw exact accuracy, the Random Forest retains a two-point advantage (0.331 vs.\ 0.310); on every metric that respects the ordinal structure of the task --- balanced accuracy, MAE, MSE, and Kendall's $\tau$ --- the additive model outperforms it. When RubricNet errs, it errs by less, and it orders pieces by difficulty more faithfully. Since the practical use case is ordinal through and through (a one-bin miss sends a student a slightly easy piece; a four-bin miss sends a beginner a concert work), the error and ranking metrics are the appropriate standard, and by that standard the interpretability of the additive model costs essentially nothing.

The absolute numbers warrant a candid reading. At 31.0\% exact accuracy over 8 classes (67.6\% over three levels, with 73.3\% of predictions within one bin), the model is not a replacement for human grading; it functions as a transparent second opinion and pre-filter over large score collections. Two considerations frame these figures. First, the failed direct 20-level experiment (Section~\ref{sec:raw20}) shows that the ceiling is partly a data problem: 716 pieces cannot populate the extremes of a 20-level scale, and even the 8-class formulation leaves classes with as few as 37 examples. Second, all grades originate from a single database, effectively a single grading tradition. The manual reviews found the grades plausible, but the model necessarily inherits any systematic bias in GuitarBurst's grading, and the reported accuracies are relative to that single standard rather than to teacher consensus at large. Part of the residual error may therefore reflect noise in the labels themselves rather than deficiencies of the model.

A further observation from Table~\ref{tab:8class_results} deserves emphasis because it generalizes beyond this thesis: every substantial performance gain came from better data and better descriptors --- repairing a corrupted extraction path, recovering rhythm, encoding demand-versus-time interactions by hand --- while the model itself never changed. For interpretable-by-construction architectures, whose representational capacity is deliberately fixed, the descriptors effectively \emph{are} the model.

\section{The Multi-Dimensional Nature of Guitar Difficulty}
\label{sec:disc_dimensions}
The third research question asked whether combining left-hand, right-hand, and global descriptor dimensions improves prediction over any single dimension alone. The ablation of Section~\ref{sec:ablation} answers affirmatively: the full set outperforms every individual dimension on every metric. The left hand is the strongest single dimension, consistent with the pedagogical view that stretches, shifts, and barre work dominate classical guitar difficulty, yet even it falls measurably short of the combination; the right-hand dimension, weak in isolation, still contributes signal the other dimensions do not carry. This mirrors, for the guitar, the finding of Ramoneda et al.~\cite{ramoneda2024combining} that piano difficulty is best modeled as a combination of separate dimensions rather than a single pooled representation, and suggests that multi-dimensionality is a property of instrumental difficulty as such rather than of any particular instrument.

\section{Interpretability}
\label{sec:disc_interpretability}
The fourth research question asked to what extent the model's descriptor contributions align with pedagogical intuition. Three strands of evidence from Section~\ref{sec:qualitative} bear on this. The influence analysis shows that the descriptors doing the most work --- piece length, shift frequency, fretboard coverage, open-string content with a negative sign --- are those a teacher would name unprompted. The monotonicity analysis shows that the influential descriptors' contributions rise (or, for \texttt{open\_string\_ratio}, fall) steadily across the eight difficulty classes, so a between-class difference decomposes into descriptor movements a teacher can verify against the score. And the triangulation against Random Forest importances and model-free rank correlations shows that this learned emphasis is not an artifact of RubricNet's own training. Taken together, these results indicate that the interpretability is substantive: the trained model can be read as a difficulty rubric grounded in guitar technique.

The interpretability has a precise cost, and the interaction-descriptor experiment of Section~\ref{sec:interaction_features} locates it. The additive constraint prevents the model from learning interactions between descriptors, and appending generic interaction products does not recover them --- such products merely introduce collinearity and degrade generalization, a failure observed across all model families tested. The only route to interaction effects that worked within the additive setting is the one taken by the V3 context-aware descriptors: encoding a specific, pedagogically motivated interaction (demand per unit time) directly into a single input. Relaxing the additive constraint into a standard neural network would likely raise raw accuracy, but at the cost of the decomposable predictions that motivate the entire approach; given the purpose of the model, the constraint is retained. It should also be noted that whether practicing teachers find the per-piece rubric decompositions convincing is an empirical question that this thesis does not answer; the evidence presented here establishes internal consistency and agreement with independent statistical measures, and an external validation with teachers remains future work.
